% FAC-EN-U0026
% French source: fac_body.tex lines 938--1015.
% Authority: PDF physical pp. 25--26; printed pp. 220--221 (no. 28 complete).
% Printed-source repair: FAC-FR-C0013.

\medskip
\phantomsection\label{fac:I:4:no:28}
\textbf{28. Double complex defined by two coverings.}
Let $\mathfrak{U}=\{U_i\}_{i\in I}$ and $\mathfrak{V}=\{V_j\}_{j\in J}$ be two coverings of $X$. If $s$ is a $p$-simplex of $S(I)$ and $s'$ is a $q$-simplex of $S(J)$, we shall denote by $U_s$ the intersection of the $U_i$ for $i\in s$ (cf. \hyperref[fac:I:3:no:18]{no.~18}), by $V_{s'}$ the intersection of the $V_j$ for $j\in s'$, by $\mathfrak{V}_s$ the covering of $U_s$ formed by the sets $\{U_s\cap V_j\}_{j\in J}$, and by $\mathfrak{U}_{s'}$ the covering of $V_{s'}$ formed by the sets $\{V_{s'}\cap U_i\}_{i\in I}$.

Define a double complex
\[
C(\mathfrak{U},\mathfrak{V};\mathcal{F})
=\sum_{p,q}C^{p,q}(\mathfrak{U},\mathfrak{V};\mathcal{F})
\]
as follows:
\[
C^{p,q}(\mathfrak{U},\mathfrak{V};\mathcal{F})
=\prod\Gamma(U_s\cap V_{s'},\mathcal{F}),
\]
where the product is taken over all pairs $(s,s')$ such that $s$ is a simplex of dimension $p$ in $S(I)$ and $s'$ is a simplex of dimension $q$ in $S(J)$.

Thus an element $f\in C^{p,q}(\mathfrak{U},\mathfrak{V};\mathcal{F})$ is a system $(f_{s,s'})$ of sections of $\mathcal{F}$ over the sets $U_s\cap V_{s'}$; equivalently, with the notation of \hyperref[fac:I:3:no:18]{no.~18}, it is a system
\[
f_{i_0\ldots i_p,j_0\ldots j_q}
\in\Gamma(U_{i_0\ldots i_p}\cap V_{j_0\ldots j_q},\mathcal{F}).
\]
We may also identify $C^{p,q}(\mathfrak{U},\mathfrak{V};\mathcal{F})$ with
\[
\prod_{s'}C^p(\mathfrak{U}_{s'},\mathcal{F}).
\]
For each $s'$ there is a coboundary operation
\[
d:C^p(\mathfrak{U}_{s'},\mathcal{F})
\longrightarrow C^{p+1}(\mathfrak{U}_{s'},\mathcal{F}),
\]
and hence a homomorphism
\[
d_{\mathfrak{U}}:
C^{p,q}(\mathfrak{U},\mathfrak{V};\mathcal{F})
\longrightarrow C^{p+1,q}(\mathfrak{U},\mathfrak{V};\mathcal{F}).
\]
Writing out the definition of $d_{\mathfrak{U}}$ gives
\[
(d_{\mathfrak{U}}f)_{i_0\ldots i_{p+1},j_0\ldots j_q}
=\sum_{k=0}^{k=p+1}(-1)^k
\rho_k\bigl(f_{i_0\ldots\hat{i}_k\ldots i_{p+1},j_0\ldots j_q}\bigr),
\]
where $\rho_k$ denotes the restriction homomorphism defined by the inclusion of\footnote{The printed target intersection retains $i_k$. The section being restricted is indexed by the simplex with $i_k$ omitted, so the larger target of the inclusion must omit $i_k$ as well. The corrected French and English restore the hat; diplomatic French preserves the print.}
\[
U_{i_0\ldots i_{p+1}}\cap V_{j_0\ldots j_q}
\quad\text{in}\quad
U_{i_0\ldots\hat{i}_k\ldots i_{p+1}}\cap V_{j_0\ldots j_q}.
\]
Similarly, define
\[
d_{\mathfrak{V}}:
C^{p,q}(\mathfrak{U},\mathfrak{V};\mathcal{F})
\longrightarrow C^{p,q+1}(\mathfrak{U},\mathfrak{V};\mathcal{F});
\]
then
\[
(d_{\mathfrak{V}}f)_{i_0\ldots i_p,j_0\ldots j_{q+1}}
=\sum_{h=0}^{h=q+1}(-1)^h
\rho_h\bigl(f_{i_0\ldots i_p,j_0\ldots\hat{j}_h\ldots j_{q+1}}\bigr).
\]
It is clear that
\[
d_{\mathfrak{U}}\circ d_{\mathfrak{U}}=0,
\qquad
d_{\mathfrak{U}}\circ d_{\mathfrak{V}}
=d_{\mathfrak{V}}\circ d_{\mathfrak{U}},
\qquad
d_{\mathfrak{V}}\circ d_{\mathfrak{V}}=0.
\]
Setting
\[
d'=d_{\mathfrak{U}},
\qquad d''=(-1)^p d_{\mathfrak{V}},
\]
endows $C(\mathfrak{U},\mathfrak{V};\mathcal{F})$ with the structure of a double complex. We may therefore apply to $K=C(\mathfrak{U},\mathfrak{V};\mathcal{F})$ the definitions of the \hyperref[fac:I:4:no:27]{preceding number}. The groups or complexes denoted in the general case by
\[
H^n(K),\quad H_I^{p,q}(K),\quad H_{II}^{p,q}(K),\quad K_I,\quad K_{II}
\]
will here be denoted respectively by
\[
H^n(\mathfrak{U},\mathfrak{V};\mathcal{F}),\quad
H_I^{p,q}(\mathfrak{U},\mathfrak{V};\mathcal{F}),\quad
H_{II}^{p,q}(\mathfrak{U},\mathfrak{V};\mathcal{F}),\quad
C_I(\mathfrak{U},\mathfrak{V};\mathcal{F}),\quad
C_{II}(\mathfrak{U},\mathfrak{V};\mathcal{F}).
\]

The definitions of $d'$ and $d''$ immediately give:

\phantomsection\label{fac:I:4:no:28:prop:2}
\textsc{Proposition 2.} $H_I^{p,q}(\mathfrak{U},\mathfrak{V};\mathcal{F})$ \textit{is isomorphic to}
\[
\prod_{s'}H^p(\mathfrak{U}_{s'},\mathcal{F}),
\]
\textit{where the product is taken over all simplices of dimension $q$ in $S(J)$. In particular,}
\[
C_{II}^q(\mathfrak{U},\mathfrak{V};\mathcal{F})
=H_I^{0,q}(\mathfrak{U},\mathfrak{V};\mathcal{F})
\]
\textit{is isomorphic to}
\[
\prod_{s'}H^0(\mathfrak{U}_{s'},\mathcal{F})
=C^q(\mathfrak{V},\mathcal{F}).
\]

We shall denote by $\iota''$ the canonical isomorphism
\[
\iota'':C(\mathfrak{V},\mathcal{F})
\longrightarrow C_{II}(\mathfrak{U},\mathfrak{V};\mathcal{F}).
\]
Thus, if $(f_{j_0\ldots j_q})$ is an element of $C^q(\mathfrak{V},\mathcal{F})$, then
\[
(\iota''f)_{i_0,j_0\ldots j_q}
=\rho_{i_0}(f_{j_0\ldots j_q}),
\]
where $\rho_{i_0}$ denotes the restriction homomorphism defined by the inclusion of
\[
U_{i_0}\cap V_{j_0\ldots j_q}
\quad\text{in}\quad
V_{j_0\ldots j_q}.
\]

Of course, the analogue of \hyperref[fac:I:4:no:28:prop:2]{Proposition~2} holds for $H_{II}^{p,q}(\mathfrak{U},\mathfrak{V};\mathcal{F})$, and there is an isomorphism
\[
\iota':C(\mathfrak{U},\mathcal{F})
\longrightarrow C_I(\mathfrak{U},\mathfrak{V};\mathcal{F}).
\]
